
\chapter{Introduction}\label{ch:paper_introduction}

\section{Motivation and Previous Work}

\paragraph{Verifiable Random Functions.} 

A verifiable random function (VRF) introduced by Micali, Rabin and Vadhan in \cite{FOCS:MicRabVad99} 
is a cryptographic primitive which allows the holder of a secret key to evaluate a pseudorandom function 
at arbitrary points and obtain image-proof pairs. In the absense of the proof the images 
are computationally indistinguishable from random elements in the range while the corresponding proof 
allows the owner of the corresponding public key to verify whether the output was indeed computed correctly.
Additionally we require that the verification accepts an image-proof pair if and only if the image is the unique 
evaluation of the point under the given secret key. As such, a VRF can be considered the public-key equivalent 
of the more common secret-key pseudorandom function (PRF).

In practice VRFs are used to implement signatures with strong verifiability guarantees \cite{C:GolOst92},
resettable zero-knowledge proofs \cite{C:MicRey01}, lottery and e-cash systems \cite{RSA:MicRiv02a,Liang2020,EPRINT:AuSusMu07a}
and in Proof-of-Stake blockchains through distributed leader election protocols  \cite{ChenMicali2016, C:KRDO17, EC:DGKR18, EPRINT:BASV23b}.

In addition to standard VRFs, both stronger variants such as constrainable VRFs which 
additionally allow the constraining of secret keys to a specified subset of evaluations \cite{},
key-homomorphic and aggregable VRFs that allow the aggregation of multiple image-proof pairs into a single
image-proof pair \cite{} and unbiasable VRFs that guarantee unbiasability, preventing adversaries 
to generated keys such that the corresponding evaluations can be biased \cite{}.

We also consider weaker notions of VRFs, such as verifiable unpredictable functions (VUFs)
already introduced in \cite{FOCS:MicRabVad99} together with VRFs and weakly verifiable random functions 
(wVRFs) which only require a weaker notion of pseudorandomness, that is pseudorandomness on uniformly random 
evaluation points. While VRFs can be constructed from VUFs using hardcore bits \cite{}, the authors of 
\cite{} show that VRFs can not, in general, be constructed from wVRFs using a seperation relying on 
enhanced trapdoor permutations.

\paragraph{Robust (Black-Box) Combiners.} 
Robust combiners are constructions that take many instances of the same cryptographic primitive, called candidate schemes,
in form of a description of their Turing machine, and securely implement that cryptographic primitive if a threshhold of the candidate schemes are secure.
Since we generally intend for the combiner to derive its security from the subset of secure candidate schemes,
and not by implementing a secure instance of the underlying primitive itself and ignoring the candidates,
we require black-box reductions from the combiner construction to any sufficiently large subset of its candidates.
In this work mainly consider robust black-box combiners, a subclass of combiners that restricts access of the combiner 
to its candidate schemes to oracles which can be queried but not inspected for their underlying description. 
Robust combiners allow the hedging of trust in cryptographic hardness assumptions underlying computational security properties,
by basing their security on a subset of hard problems rather than a a single hard problem. 
In practice robust combiners are uesful to reduce the dependency of a systems 
on the security of a specific third party implementation, e.g. a cryptographic library.
In theory robust combiners are related to universal cryptographic construction, i.e.
constructions that implement a secure instance of a cryptographic primitive under the only assumption 
that a secure instance of the cryptographic primitive exists, a prominent example is Levin's universal one-way function\cite{}.
Finally combiners play a central role in the transition towards post-quantum cryptography where 
hardness assumptions are not as well studied as the respective assumptions for the classic cryptography counterpart, 
in particular the combination of classical and post-quantum cryptographic assumptions has been a fruitful research field. 

Combiners have been considered for a wide variety of cryptographic primitivey, including 
key exchange protocols \cite{}, signature schemes \cite{}, oblivious-transfer \cite{}, zero-knowledge proofs \cite{}, 
commitment schemes \cite{}, and indistinguishability obfuscation \cite{}. While pseudorandom functions (PRFs) 
admit easy robust $(n;1)$ combiners by simply XOR'ing the outputs of the candidate schemes, 
evaluated on independent keys the combination of VRFs, the combination of VRFs is less obvious.
To the best of our knowledge, there has been no prior research on robust combiners for verifiable random functions.

\section*{Our Contribution}

\paragraph{Robust Black-Box DAG Combiner.}
We introduce the notion of robust black-box DAG combiners, formally defined in 
\Cref{def:robust_blackbox_dag_combiner}, DAG combiners for short.
DAG combiners are black-box combiners which are defined by directed acyclic graphs (DAGs) for every 
implemented algorithm. The implemented algorithms are specified through that DAG and their execution 
follows the DAG from a source vertex to a target vertex, where all intermediate vertices correspond to 
an oracle query to a specified candidate on a valid computed from the output of 
all predecessor outputs. DAG combiners make the structure of the black-box combiner more explicit 
by structuring it into interleaved oracle queries to the candidats and intermediate computation,
that we model with Turing machines without oracle access to the candidates. 
To connect DAG combiners and general black-box combiners we show that the notions are indeed equivalent,
that is, that for every general black-box combiner we can construct a DAG combiner 
that behaves identically in the sense that the algorithms defined by the combiners are equivalent.
We state this result formally in \Cref{thm:generality_dag_comb}.

\paragraph{Robust Black-Box VRF Combiner.}
When we look for robust black-box combiner constructions we typically consider 
$(2;1)$ combiners first because they naturally give rise to $(n;1)$ combiners by 
iteratively applying them. We prove that there exists no robust $(2;1)$ black-box combiner for the 
VRF primitive by showing that against any ppt. Turing machine with access to two 
candidate schemes, one of which implements a secure VRF, that implements a 
uniquely provable and perfectly correct verifiable function there exists an adversary with access
to a PSPACE complete oracle that wins the PRF security experiment against this construction 
with high advantage.
The impossibility of a robust $(2;1)$ black-box VRF-combiner implies that there exists no robust
$(n;\lfloor n/2\rfloor)$ black-box combiner for VRFs. Given $n$ oracles, this means that there is no 
way to combine them a way such that the construction guarantees VRF security without the assumption 
that more than half the candidates are secure.

Given the impossibility of $(n;\lfloor n/2\rfloor)$ black-box VRF combiners we ask whether such combiners 
exist if we strengten the assumptions on the candidate schemes or weaken the requirements on the combiner.
One idea is that the information theoretic properties on the combiner post too strong a requirement such
that the constructions are not possible, that is however not the case. We show that the impossibility 
holds even if we lower the requirements to statistical and computational notions of correctness and unique provability.
Further we show that even pseudorandomness can be relaxed, not changing the feasibility.
Statistical notions of unique provability and perfect correctness are rarely considered in literature,
the impossibility of a combiner under those assumptions implies that we have to strengthen the assumptions 
on the combiners. Finally we show that even if we require all candidates to be pseudorandom, that there exists no 
combiner construction. 

On the other hand if require all candidates to be unqiuely provable and perfectly correct then very 
simple constructions are possible that are based on the XOR PRF-combiner. Such a combiner implements 
its key generation algorithm by independently sampling a key pair for every candidate and then 
combining all secret keys and all public keys in a tuple, forming the combiner keys.
The evaluation algorithm evaluates a point $x$ on all candidate schemes under their respective 
secret key to obtain image-proof pairs $(y_i,\pi_i)$ for every candidate. The final combiner output 
is given by the xor of the candidates outputs $y = y_1 \oplus y_2$ while the proof is given by all 
computed intermediate results $\pi = (y_1, y_2, \pi_1, \pi_2)$. The verification algorithm works by 
querying both candidates to check that $(y_1, \pi_1)$ and $(y_2, \pi_2)$ are accepted by the individual oracles,
then it checks that $y = y_1 \oplus y_2$ and accepts iff both hold. The resulting $(2;1)$ weak black-box combiner 
for VRFs can be used to construct a $(n;1)$ weak black-box combiner for VRFs.

Having shown that weak $(n;1)$ black-box combiner exist we study weak combiners for more advanced VRF notions.
Specifically we inspect the possibility of combining constrainable VRFs, key-homomorphic VRFs, aggregable VRFs 
and unbiasable VRFs. \todo{expand if time}

\begin{comment}
We introduce a new definition of a robust black-box combiner, a robust black-box DAG combiner,
which models the implemented algorithms of the combiner as computations that follows the structure 
of a vertex-colored directed acyclic graph starting from a source and ending in a sink, where for each vertex 
in the graph we query a specific oracle of a specific candidate on the output of a probabilistic polynomial 
time Turing machine of all predecessors. We show that the notion of a robust black-box DAG combiner 
is sufficiently strong to capture all ppt. robust black-box combiners. We use this new model to show 
the impossibility of combining multiple instances of wVRFs which extends to the impossibility of combining 
VUFs and VRFs without additional assumptions. We argue that the combination of the unique provability property 
is however less important since it is of information theoretic or statistical nature and therefore 
doesn't rely on any computational hardness assumption. In fact, good VRF constructions such as ECVRF \cite{},
lattice VRF \cite{} and \cite{} mathematically prove unique provability mathematically, rendering 
combination less important. We consider the possibility to combine more advanced security properties 
on VRFs such as unbiasability for which we construct a combiner for uniform constructions which relies on collision resistant
hash functions and randomness extractors to combine unbiasable VRFs under the additional assumption that 
all candidate schemnes are uniquely provable. We also investigate combiners for secret-key aggregability and public-key 
aggregability. Secret-key aggregability is a publicly verifiable property due to an efficient algorithm by 
Blum, Luby and Rubinfeld \cite{} which allows checks for the homomorphism property of a function 
in time independent of the size of the underlying groups. Combiners for publiy-key aggregability 
\todo[]. Finally we investigate additional properties such as constrainability \cite{} \todo[] and which all admit the 
simple pseudorandomness combiner construction.
\end{comment}


% \section*{Thesis Outline}
% 
% \begin{itemize}
%     \item \textbf{\Cref{ch:notation} (Notation and Preliminaries):}
%     \item \textbf{\Cref{ch:analysis} (Algorithm and Analysis):} 
%     \item \textbf{\Cref{ch:conclusion} (Conclusion):} 
% \end{itemize}
